% ---
% id: protocol-compiler/def:validator
% kind: def
% title: "The emitted validator"
% summary: >-
%   Eight exact checks on a candidate declaration: sector maps, the marginal,
%   the operators, the key alphabet, sifting, constraint typing, feasibility
%   and party computability. Feasibility passes only on an empty list, a
%   verified default candidate or a verified supplied witness.
% prominence: standard
% status: implemented
% uses: [def:declaration, def:exact, thm:feasibility, def:compile, def:spec, def:homog, def:target]
% mentions: [thm:validator-soundness, def:chaincond, thm:measurements-and-statistics, thm:sifting-and-key, rem:chain-and-validator, req:validator-refusals, req:derived-objects, foundations/req:exact-island]
% assumes: []
% refs: []
% provenance: newly-derived
% implements: certiqs_key.protocol_compiler.validation:validate
% lean:
% notes: "thm:measurements-and-statistics is a mention and not a use: V1's pointer to it explains why V1 has no off-block conjunct and no check rests on it, which lets thm:typing and thm:measurements-and-statistics use this definition without a cycle."
% ---

The validator is the list of checks the compiler runs on a candidate
declaration before it emits one. A candidate failing a check is refused with a
reason naming the check (Requirement~\ref{req:validator-refusals}); what
passing all eight establishes is Theorem~\ref{thm:validator-soundness}, and which conditions
are decided on the chain rather than here, and which objects check V1 does
not compare, is Remark~\ref{rem:chain-and-validator}. That check V7 decides
on quantities formed in exact arithmetic is
\xref{foundations/req:exact-island} at work, and that no product block,
\tref{lifted operator} or denoted family of an \tref{ensemble form} is a
datum is
Requirement~\ref{req:derived-objects}.

\begin{definition}\label{def:validator}
Given a candidate declaration \eqref{eq:declaration} with its emitted
\tref{index map}, the checks are as follows. \emph{Every operator condition
below is decided exactly}, in the sense of Definition~\ref{def:exact}, and
every combinatorial condition exactly as well.
\begin{description}[leftmargin=3.0em,style=nextline]
\item[V1 (finiteness, totality, and the \tref{sector map})] every declared alphabet is
  finite and non-empty and every declared map is tabulated on all of its stated
  domain; every declared system carries a finite dimension;
  $d_{\Qsys}=\sum_{s\in\SecSet}d_s$; the restriction of $\idxm$ to $\SecSet$ is
  the total order the sector map fixes; \emph{and the sector maps
  agree}: there is one emitted map $\Smap$, and the map carried by each
  supplied $\Qsys$-operator family is compared against it in \emph{block
  kind}, in the \emph{list of block identifiers}, in their \emph{order} and in
  their \emph{dimensions}. A disagreement on any of the four is refused,
  naming the first position at which the lists differ. The families compared
  are exactly the $\Qsys$-operator families of \eqref{eq:declaration}, which
  is the list D(i) quantifies over and no wider one: the emitted family
  $\{F_v\}$, which the last link of $\Chrec$ delivers as its \tref[target povm]{target POVM}
  against the map the chain's producer wrote it against, and each constraint
  entry of $\CList$, declared in $\Spec_{\mathrm{con}}$ or transported by a
  link of either segment. Those two slots are the whole of
  \eqref{eq:declaration} that carries a $\Qsys$-operator
  (Definition~\ref{def:declaration}), so the list is exhaustive by
  construction rather than by enumeration; the objects outside it are compared
  where they belong. The entries of
  $\Spec_{\mathrm{con}}$ are compared as well as the transported ones:
  Definition~\ref{def:compile} fills $\CList$ from both sources,
  Definition~\ref{def:spec}\,(e) has each declared $\hat O_k$ carry its own
  sector map, and an entry written against a different map is the error this
  check exists for, whichever source it arrived from. The comparison is
  combinatorial --- it compares labels and integers, never matrix entries ---
  and it is the whole of the structural verification the emitted declaration
  receives: block-diagonality is carried by the type
  (Theorem~\ref{thm:measurements-and-statistics}\,(ii)), so the one structural error that can still
  arise is two parties disagreeing about what the sectors are, and V1 is the
  check for it.
\item[V2 (\tref{marginal slot})] the slot holds $\nomarg$, or a matrix $\tau_A$ on $A$
  --- one block, $A$ carrying no sector decomposition --- satisfying
  $\tau_A=\tau_A^{\dagger}$, $\tau_A\succeq0$ and $\tr\tau_A=1$, each decided
  exactly, and with the trace accumulated over the row index verified to
  \tref[rounds nowhere]{round nowhere}; a failure of the last is refused.
\item[V3 (operators, block by block)]
  $M_x=M_x^{\dagger}$ and $M_x\succeq0$ for each $x\in\Xset$, and
  $\sum_xM_x=\idon{A}$ with the sum, in the emitted index order, verified to
  round nowhere, a failure of which is refused;
  $F_v^{(s)}=F_v^{(s)\dagger}$ and $F_v^{(s)}\succeq0$ for each
  $(v,s)\in\Vset\times\SecSet$, and, for each $s$,
  $\sum_vF_v^{(s)}=\idon{\Qsys_s}$ with that sum verified to round nowhere, a
  failure being refused with the sector and the outcome index at which the
  first inexact addition occurred; and
  $\hat O_k^{(s)}=\hat O_k^{(s)\dagger}$ for each $(k,s)$. Every one of these
  is decided exactly. There is no off-block condition: an emitted
  $\Qsys$-operator is a \tref{block list}, so there is no $\Pi_sF_v\Pi_{s'}$ or
  $\Pi_s\hat O_k\Pi_{s'}$ to evaluate and nothing that could be non-zero.
\item[V4 (key alphabet)] every tabulated value of $r$ lies in
  $\Zval\sqcup\{\perp\}$.
\item[V5 (sifting consistency)] for all $(x,v)$ and $(x',v')$ with
  $a(x,v)=a(x',v')$: $r(x,v)=\perp\iff r(x',v')=\perp$.
\item[V6 (constraint typing, including the tag)] when the marginal
  slot is filled: $\bowtie_k\in\{\preceq,=,\succeq\}$, and the entry's
  right-hand side is of one of the two admitted types, which is what fixes its
  tag (Definition~\ref{def:homog}): \tref{homogeneous}, whose datum is of
  real scalar type, or \tref{general}, whose datum is a matrix on $A$ with
  $\Theta_k=\Theta_k^{\dagger}$ decided exactly. When the slot holds
  $\nomarg$: $\bowtie_k\in\{\le,=,\ge\}$ and the right-hand side is of real
  scalar type. A real scalar right-hand side is accepted in
  \tref{trusted-source mode} at every $d_A$, and is the licensed form there
  (Definition~\ref{def:homog}). The relation symbol, the type of the
  right-hand side and the tag that type fixes are compared exactly; an entry
  whose right-hand-side slot holds neither admitted object, or a complex-typed
  right-hand side in \tref{untrusted-source mode}, is refused by type and not
  measured. The check also reads an object \eqref{eq:declaration} does not
  carry, the tag field of an \emph{arriving} entry, in $\Spec_{\mathrm{con}}$
  or transported by a link: where an arriving entry carries an explicit tag
  field, V6 compares that field with the type of the datum it names, and a
  disagreement is refused. An entry of \eqref{eq:declaration} in
  either of the two forms above is the triple
  $(\hat O_k,\bowtie_k,\mathrm{rhs}_k)$ and holds no tag field independent of
  the type of its datum, the emitted tag being read off that type
  (Definition~\ref{def:homog}); so on a candidate declaration this clause has
  no subject, refuses nothing, and is not among the conjuncts D(vi) is paired
  with (Theorem~\ref{thm:validator-soundness}, proof). An entry carrying no tag field is not
  refused: its tag is the type of its datum and is read off it, which is what
  lets an entry declared as
  a real constant, carrying no tag field, be read here.
\item[V7 (feasibility; one-sided, three branches)] V7 passes in exactly three
  cases:
  \begin{enumerate}[label=(\alph*),leftmargin=2.2em]
  \item $\CList=\varnothing$; or
  \item $\CList\ne\varnothing$ and the default candidate
    \eqref{eq:witness-empty}, which the compiler forms and does not emit,
    passes the per-block verification of Theorem~\ref{thm:feasibility}\,(ii); or
  \item $\CList\ne\varnothing$, the default candidate fails, and a witness
    $\omega^{\star}$ is supplied \emph{as a block list against $\Smap$} --- the
    map it carries compared with $\Smap$ in block kind, in identifiers, in
    their order and in their dimensions, exactly as V1 compares the map of a
    declared $\Qsys$-operator family --- and passes that same verification.
    This is the one place a witness's sector map is compared, and it is
    here rather than at V1 because $\omega^{\star}$ is a block list on
    $A\Qsys$, is no slot of \eqref{eq:declaration}, and is quantified over by
    no target condition (Definition~\ref{def:target}).
  \end{enumerate}
  Otherwise V7 fails and the candidate is refused, the refusal naming, where a
  witness was supplied and rejected on its map, the first position at which
  the two lists differ. The verification in (b) and (c) is the $|\SecSet|$
  block positivity decisions --- Hermiticity and semidefiniteness per block
  --- the trace, the marginal equality in trusted-source mode, and one leg per
  constraint entry, all decided exactly, on quantities formed in exact
  arithmetic (Theorem~\ref{thm:feasibility}\,(ii)). It is taken on a
  candidate satisfying D(i)--D(iii) and D(vi), which V1--V3 and V6 decide; a
  candidate failing one of them is refused there. V7 is the one check of the
  list that
  is \emph{not} equivalent to its target condition: it implies D(vii) and is
  implied by nothing weaker than a candidate that verifies
  (Theorem~\ref{thm:validator-soundness}\,(iii)).
\item[V8 (party computability)] $a_A$ and $a_B$ are tabulated on
  $\Oset{\mathsf A}$, $\Oset{\mathsf B}$; and for all $(x,v)$, $(x',v')$ with
  $\rd{\hold}(x,v)=\rd{\hold}(x',v')$ and $a(x,v)=a(x',v')$:
  $r(x,v)=r(x',v')$.
\end{description}

\emph{The general format.} On a candidate declaration in the general format
of Definition~\ref{def:declaration}\,(DX1)--(DX6) the eight checks read as
follows, every operator condition again being decided exactly. In the
\tref{basic format} --- one Alice sector, one \tref{measured factor},
\tref{POVM form}, two messages neither reading the other and no party-local
entry --- each reading is the check of the
same name above, conjunct for conjunct, so on such a candidate the two
readings admit and refuse alike.

\emph{V1 in the general format.} Finiteness and totality as above, every
table of the transcript included: each $a_m$ is tabulated on all of
$\Oset{\pi_m}\times\mathcal{C}_1\times\dots\times\mathcal{C}_{m-1}$, which is
also where the signature clause of D(viii) lives: a table whose stated
domain names a later message's alphabet is not of that type and is refused
here. \emph{The sector map is one factorised map}: at most one factor tagged
$A$ and, where there is one, first in the list; at least one factor tagged
$\Qsys$; \emph{one block kind across the factors}, a disagreement being
refused, naming the first factor at which the kind differs;
$d_{\Qsys_f}=\sum_{s_f}d_{f,s_f}$ for every factor and
$d_{\Qsys}=\prod_fd_{\Qsys_f}$; and, where an $A$ factor is declared,
$d_A=\sum_td_{A_t}$. The restriction of $\idxm$ to each factor's \tref{sector set}
is the total order that factor's own map fixes, and its restriction to
$\mathbb{T}$ is the order the $A$ factor's map fixes; the product orders on
$\SecSet$ and $\Vset$ and the order on the \tref[joint sector]{joint sectors} are derived
by (DX1) and (DX3) and are not compared, there being no second datum to
compare them with. \emph{And the maps agree, factor by factor}: the map
carried by each declared family $\{F^{f}_{v_f}\}_{v_f}$ is compared against
the map of factor $f$, and the map carried by each entry's operator
$\hat O_k$ against the map of that entry's scope, in the same four
components, with the same refusal naming the first differing position.
\emph{Alice's \tref{declared data} are compared too, against the map of the
factor tagged $A$}: the map carried by $\tau_A$, the map carried by each
operator of Alice's measurement slot --- $M_x$ in POVM form, $N_x$ in
\tref{ensemble form} --- and the map carried by each $\Theta_k$ of an entry declared
general are each compared against that factor's map in the same four
components (block kind, block identifiers, their order, dimensions), with the
same refusal naming the first differing position. That comparison is what
makes condition~(T) of Definition~\ref{def:exact} decided for Alice's data,
which its general-format paragraph calls non-empty for them; it is decided
here and nowhere else, and without it a candidate whose Alice block list
carried a different block kind, a different identifier list or a different
order would pass every check and violate D(i). Where no factor is tagged $A$
there is one Alice block and no map to compare, as in V1 above. \emph{The list
of compared objects is again exhaustive by construction}: the objects D(i)
quantifies over are the per-factor families, the entries' operators and
Alice's declared data. The product family is no datum, and the
\tref{lifted operator} of a scoped entry is no datum; the descriptors that the
last links of a chain produce, decided per factor at
Definition~\ref{def:chaincond}\,(C1), and a supplied witness, compared at
branch~(c) of V7, stay outside the list; and the physical families stay outside
it too, one per factor.

\emph{V2 in the general format.} The slot holds $\nomarg$, or a block list
$(\tau^{(t)}_A)_t$ against the $A$ factor's map with
$\tau^{(t)}_A=\tau^{(t)\dagger}_A$ and $\tau^{(t)}_A\succeq0$ for every $t$
and $\sum_t\tr\tau^{(t)}_A=1$, each trace accumulated over the row index and
the block traces then added in the \tref{declared order} of $\mathbb{T}$, with
\emph{every} one of those additions verified to round nowhere; a failure of
the last is refused. \emph{Where Alice's measurement slot is in ensemble form,
additionally $\tau^{(t)}_A\ne0$ for every $t$}, a comparison of declared
entries with the datum $0$, a failure being refused naming the sector. At
$|\mathbb{T}|=1$ this is V2 above, the one block being the whole of $A$.

\emph{V3 in the general format.} \emph{Per measured factor}: for every
$f\in\mathbb{F}$, every $v_f\in\Vset^{f}$ and every $s_f\in\SecSet^{f}$,
$F^{f,(s_f)}_{v_f}$ Hermitian and positive semidefinite, and for each $s_f$
the sum $\sum_{v_f}F^{f,(s_f)}_{v_f}=\idon{\Qsys_{f,s_f}}$ with its
declared-order sum verified to round nowhere, a failure being refused with
the \emph{factor}, the sector and the outcome index at which the first
inexact addition occurred. \emph{On Alice's data, block by block along
$\mathbb{T}$}: in \tref{POVM form} $M^{(t)}_x$ Hermitian and positive
semidefinite for every $(x,t)$ and $\sum_xM^{(t)}_x=\idon{A_t}$ for every
$t$, the sum in the emitted index order verified to round nowhere, a failure
being refused with the Alice sector; in \tref{ensemble form} $N^{(t)}_x$
Hermitian and positive semidefinite for every $(x,t)$ and, \emph{in place of
the completeness identity}, the identity $\sum_xN^{(t)}_x=\tau^{(t)}_A$ of
\eqref{eq:ensembleform} decided entry by entry against the declared block of
$\tau_A$, a failure being refused with the sector and the first differing
entry, with that declared-order sum verified to round nowhere, a failure
being refused with the sector and the index of the first inexact addition.
\emph{An ensemble-form slot beside a marginal slot holding $\nomarg$ is
refused}; this is D(iii)'s general-format clause, and it is a decision on
which of two slots is filled, not on a number. And $\hat O_k^{(s_{f_k})}$ is
Hermitian for every entry and every block of its scope. \emph{Nothing is
decided on a product block \eqref{eq:productfamily}, on a lifted operator or
on the denoted family \eqref{eq:ensemblepovm}}: none of the three is a datum,
and that the first is a POVM on
$\Qsys$ and the third a POVM on $A$ are theorems
(Definition~\ref{def:declaration}\,(DX3), Theorem~\ref{thm:sifting-and-key}\,(v)).

\emph{V4, V5, V7 in the general format.} V4 and V5 keep their wording, with
$a$ the function \eqref{eq:transcript} derives from the transcript's tables;
V5's pairs are compared on the values of that function and on nothing about
how it was built. V7 keeps its three branches and its refusal, and what its
branches (b) and (c) verify is the per-\emph{joint}-sector verification of
Theorem~\ref{thm:feasibility}\,(ii) in the general format, on a candidate
satisfying D(i)--D(iii) and D(vi), in their general-format reading, which
V1--V3 and V6 decide (a candidate failing one of them is refused there): the
default candidate is
\eqref{eq:witness-empty} read per joint sector, and a supplied
$\omega^{\star}$ is required to be a block list over the joint sectors, its
map compared in the same four components at branch~(c) and there alone.

\emph{V6 in the general format.} When the marginal slot is filled:
$\bowtie_k\in\{\preceq,=,\succeq\}$, and the entry's right-hand side is of
one of the \emph{three} admitted forms, which with the arriving tag field is
what fixes its tag (Definition~\ref{def:homog}): homogeneous with a real
scalar, general with a block list $(\Theta^{(t)}_k)_t$ along $\mathbb{T}$
decided Hermitian block by block, or \tref{party-local} with a real scalar. When the
slot holds $\nomarg$: the right-hand side is of real scalar type, either
untagged or tagged party-local. \emph{The relation symbol is decided by form
and not by mode alone}: for an \emph{untagged} entry
$\bowtie_k\in\{\preceq,=,\succeq\}$ when the slot is filled and
$\bowtie_k\in\{\le,=,\ge\}$ when it holds $\nomarg$, exactly as in V6 above;
for an entry declared party-local, $\bowtie_k\in\{\preceq,=,\succeq\}$
\emph{in either case}, its reading being an operator relation
(Definition~\ref{def:homog}, the relation-symbol clause of the \tref[party-local form]{party-local
form}; Definition~\ref{def:target}, D(vi) in the general format), with the
scalar symbols admitted as well where every row of that entry is $1\times1$
--- every Alice block, and every sector block of every measured factor other
than the entry's scope, of dimension one, $d_A=1$ at one measured factor
included --- the three operator symbols being the three scalar symbols there.
This is a decision on a declared symbol beside a declared tag and measures
nothing, and at an entry not declared party-local it is the clause of V6
above word for word, so nothing that clause admits or refuses moves.
\emph{Every entry's scope names a measured factor of the emitted sector map},
a comparison of a label against the factor list, a failure being refused
naming the entry; this is the general-format half of D(vi) and is decided and
never measured, as the relation symbol is. \emph{A party-local entry is the
one that arrives carrying that tag value}, and it carries a tag field the
type of its datum does not determine, that type not separating it from a
homogeneous one. On a candidate carrying one, the arriving-tag clause of V6
above therefore has a subject, and it is this comparison: the datum is of the
type \emph{that} tag admits, a real scalar, and the scope names a measured
factor; a party-local tag beside a matrix right-hand side is refused. The
sentence of V6 above that an entry holds no tag field independent of the type
of its datum is true of the two untagged forms, homogeneous and general, and
of them alone. \emph{Nothing is inferred}: an entry arriving without the
party-local tag is never read party-locally. So the arriving-tag clause
refuses no party-local entry beside a real scalar, which a literal reading of it
would otherwise do, and wherever no entry is tagged the reading is that of V6
above (Definition~\ref{def:homog}).

\emph{V8 in the general format.} The announcement clause reads the
transcript: every $a_m$ is tabulated on all of its stated domain, which is
V1's clause restated where D(viii) pairs with it; the key clause is unchanged,
with $a$ the derived function.
\end{definition}
